\chapter{Поверхностное натяжение в двухфазных двухкомпонентных системах}\label{ch:ch3}

\section{Использование молекулярной динамики для моделирования фазовых переходов}

Молекулярная динамика позволяет применить множество методов и рассчитать некоторые интересные величины, в том числе недоступные для эксперимента.

Самый очевидный способ моделирования фазового перехода первого рода \--- создание системы в одном из неравновесных состояний, в котором произойдет фазовый переход в результате установления в системе равновесия.  Этого можно добиться охлаждением или сжатием пара для того, чтобы добиться конденсации. В результате этого можно определить такие величины как размер критического зародыша или скорость нуклеации. Интерес представляет и форма зародыша, которая может быть довольно нетривиальна. Недостаток этого способа заключается в невозможности определить температуру фазового перехода из-за наличия метастабильного состояния.

Развитием этого метода является Z-метод~\cite{Belonoshko2006}. Поскольку моделирование происходит при постоянных объёме, температуре и энергии, он лучше подходит для фазового перехода кристалл-жидкость и других, в которых плотности фаз сопоставимы. Он заключается в моделировании при постоянной энергии систем с различной начальной температурой. Система приходит в равновесное состояние с тремя возможными исходами. Либо начальная энергия окажется недостаточной, чтобы достичь температуры плавления, и произойдет только нагрев. Если энергия окажется слишком большой, произойдет полное расплавление системы и дальнейшее нагревание. В промежуточном случае между этими двумя установится двухфазное равновесие при некоторой температуре, которая и будет температурой плавления. 

В случае, если разность плотностей фаз достаточно велика, как в случае жидкой и газовой фаз, можно использовать двухфазное моделирование. Этот метод включает в себя создание двухфазной системы с плотностями, близкими к плотностям фаз при данных условиях и вывод такой системы на равновесие. Таким образом можно получить равновесные плотности и поверхностное натяжение. Этот метод был использован в данной работе для моделирования смеси метана и этана. 

В подразделе \cref{subsec:ch3/sect2/sub1} описан процесс моделирования, включающий использованную модель, начальное состояние системы, процесс выхода на равновесие и метод расчета поверхностного натяжения. Подраздел \cref{subsec:ch3/sect2/sub2} посвящен расчету поверхностного натяжения чистого этана и предсказанию его значений с помощью парахора. В ~\cref{subsec:ch3/sect2/sub3} при ведены результаты расчета поверхностного натяжения многокомпонентной системы. Результаты, приведенные в данной главе, опубликованы в~\cite{RJOPC285654}.

\section{Поверхностное натяжение}

Для расчета поверхностного натяжения методом двухфазного моделирования характерны некоторые сложности. Некоторые из них \--- долгий выход системы на равновесие и необходимость набора статистики в течение длительного времени. Это требует большого количества вычислительных ресурсов. Чтобы сократить затраты, можно использовать более простую модель. Однако, из-за этого пострадает точность. 

Ещё одной проблемой является невозможность использовать дальнодействующие поправки при двухфазном моделировании. Силовые поля параметризуются таким образом, чтобы при включении этих поправок давать корректные плотности. При отключении их разность экспериментальной и полученной методом молекулярной динамики плотностей будет составлять до 5 процентов. Однако, при двухфазном моделировании включение поправок приведет к некорректному описанию границы раздела фаз в силу неоднородности системы в её окрестностях. И именно поведение вещества в этой области влияет на величину поверхностного натяжения.

Следует заметить, что поверхностное натяжение является крайне чувствительной к точности расчета характеристикой: её изменение пропорционально 4-ой степени разности плотностей жидкости и газа. Потому даже небольшое изменение плотности, например, на 5 процентов, приведет к изменению конечного результата на около 20 процентов. 

Несмотря на это, расчет поверхностного натяжения методом молекулярной динамики можно проводить, и при этом даже не требуется использовать модель, которая идеально воспроизводит плотности жидкой и газовой фазы. Молекулярная динамика хорошо воспроизводит зависимость поверхностного натяжения от плотности. Благодаря этому свойству можно ввести поправки к плотности и получить корректное значение поверхностного натяжения


\subsection{Расчет поверхностного натяжения методом молекулярной динамики}\label{subsec:ch3/sect2/sub1}

В этой части работы используется модель Transferable Potential for Phase Equilibria -- United Atom (TraPPE-UA) \cite{trappeua}. Атомы углерода объединяются со связанными с ними атомами водорода в единые псевдоатомы, что позволяет значительно ускорить расчеты за счёт уменьшения числа взаимодействующих атомов. В данном случае это превращает молекулу метана в единый псевдоатом, а молекулу этана \--- в два псевдоатома, связанных  друг с другом.

Парные межмолекулярные взаимодействия описываются потенциалом Леннард-Джонса с радиусом обрезания 14~\AA . Для метана его параметры равны $\epsilon_{CH_4} / k_B = 148~K$, $\sigma_{CH_4} = 3.73$~\AA, для этана \--- $\epsilon_{C_2 H_6} / k_B = 98~K$, $\sigma_{C_2 H_6} = 3.75$~\AA. Перекрестные коэффициенты вычисляются с помощью правил Лоренца-Бертло. Радиус обрезания потенциала равен 17~\AA. Потенциал непрерывно соединяется с нулем через полиномиальную функцию. Её коэффициенты выбираются так, чтобы значения функции и производной слева и справа были равны на расстоянии 15~\AA, где полином соединяется с потенциалом Леннард-Джонса, и равны нулю на расстоянии 17~\AA. Этот метод позволяет отчасти компенсировать отсутствие дальнодействующих поправок и провести непрерывный переход потенциала взаимодействия к нулю.

Связь C-C в этане описывается гармоническим потенциалом:
\begin{align}
	U_{bond} = k_{C-C} (r-r_0)^2 ,
\end{align}

 с константой жесткости $k_{C-C}N_A = 268 кДж/($моль \AA $^2) $, где $N_A$ \--- постоянная Авогадро и $r_0 = 1.54$ \AA.

Такая модель хорошо описывает фазовые равновесия несмотря на описанные выше упрощения \cite{trappeua}. В то же время она значительно, примерно в пять раз, уменьшает время расчета. Это важно, поскольку при высоких температурах и при приближении к критической точке возрастает время выхода на равновесие. 

Создается длинная вычислительная ячейка со сторонами, подчиняющимися соотношению $l_x:l_y:l_z$ = 10:1:1~\cref{fig:sys_surf_tension}. В её центре, перпендикулярно оси \textit{x}, создается жидкая пленка. Для этого в соответствующую область помещаются молекулы метана и этана с выбранными концентрациями и с плотностью, близкой к плотности жидкости. В оставшейся части вычислительной ячейки располагается ещё 2000 молекул метана и этана. Их плотность будет достатчоно невелика, чтобы эти молекулы были газом. Мольная доля метана в системе варьировалась от 0 до 0.4. 

\begin{figure}[ht]
    \centerfloat{
        \includegraphics[scale=0.75]{profile_ovito.png}
    \caption{Смесь метана и этана при температуре 213 К и давлении 22 атм. Красным показаны псевдоатомы метана, синим \--- молекулы этана.}
    }\label{fig:sys_surf_tension}
\end{figure}

Далее происходит вывод системы на равновесие. Температура поддерживается постоянной на протяжении всего времени моделирования при помощи термостата Нозе-Гувера \cite{HOOVER1996253}. В течение 2~нс система приводится к нужному давлению с помощью баростата Шиноды \cite{Shinoda}. После этого моделирование идет в течение 20~нс при постоянном объёме. За это время плотности жидкой и газовой фаз становятся постоянными, что служит критерием наступления равновесия в системе. Этот этап моделирования может занимать до 10~нс в случае рассмотренных условий, близких к критической точке. При дальнейшем повышении температуры время установления равновесия возрастает ещё сильнее вплоть до того, что равновесия достичь не удалось. Оставшееся время происходит набор статистики для усредненных профилей плотности и диагональных компонент тензора давлений.

Поверхностное натяжение в молекулярной динамике зачастую~\cite{simmons_molecular_2004, ghoufi_computer_2016, ghoufi_calculation_2019, malfreyt_calculation_2014} рассчитывается через интеграл от нормальной и тангенциальной к поверхности компонент тензора давлений. В случае плоской поверхности, расположенной в $x = 0$, формула примет вид:

\begin{align}
	\sigma = \int_{-\infty}^{+\infty} (P_{xx} - (P_{yy} + P_{zz})/2) dx,
\end{align}


где $\sigma$ \--- поверхностное натяжение, $P_{xx}$, $P_{yy}$ и $P_{zz}$ \--- диагональные компоненты тензора давлений. $P_{xx}$ \--- компонента, ортогональная 2 исследуемым поверхностям, а $(P_{yy}+P_{zz})/2$ \--- параллельная им.

Поскольку подынтегральное выражение отлично от нуля только в некоторых небольших частях системы, а моделируемая ячейка \--- конечна, можно провести интегрирование и получить следующий результат:

\begin{align}
	\sigma = \frac{L_x}{2} (P_{xx} - \frac{P_{yy}+P_{zz}}{2}),
\end{align}

где \--- $L_x$ длина вычислительной ячейки в направлении оси $x$, а множитель 1/2 объясняется наличием двух межфазных границ в ячейке.

\subsection{Верификация модели расчетом парахора чистого этана}\label{subsec:ch3/sect2/sub2}

Расчет поверхностного натяжения проводился в температурном диапазоне 203-253~K. В силу правила фаз Гиббса невозможно зафиксировать одновременно температуру или давление и плотность, поскольку у системы только одна степень свободы. Потому точки получены при разных температурах, и с ростом температуры и приближением к критической точке убывает разность плотностей жидкой и газовой фаз.

С уменьшением разности плотностей падает и разность между нормальной и тангенциальной к исследуемым межфазным поверхностям составляющими тензора давлений, а с ними убывает и поверхностное натяжение. Поскольку при этом слабо меняется величина флуктуаций давления, возрастает относительная погрешность. 

Для аппроксимации зависимости поверхностного натяжения чистого этана от разности плотностей жидкой и газовой фаз использовалась эмпирическая величина парахор~\cite{_-_1996}:
\begin{equation}
    \label{eq:parachor_definition}
    \begin{alignedat}{2}
	\sigma = [P_r (\rho_{liq} - \rho_{gas})]^4
	\end{alignedat}
\end{equation}

где $\rho_{liq}$ и $\rho_{gas}$ \--- плотности жидкой и газовой фазы соответственно, $P_r$ \--- парахор~\cite{escobedo}:

$$ P_r = P_0 (1-T_r)^{0.37} T_r exp(\frac{0.30066}{T_r}+0.86442T_r^9), $$

где $ T_r=T/T_c $ \--- отношение температуры к критической температуре, $P_0$ \--- табличный параметр, зависящий от вещества.

Проведены расчеты поверхностного натяжения чистого этана при условиях, подобным выбранным в эксперименте \cite{baidakov}. Как можно заметить по графику, при одних и тех же температурах в системе не совпадают разности плотностей полученные в эксперименте и методом молекулярной динамики. Так выражает себя недостаточный учет дальнодействия в модели. Энергия взаимодействия получается ниже фактической, меньше силы притяжения между молекулами и ниже плотность. 

Следует заметить, что хотя по оси абсцисс и отложены разности плотностей, величина плотности жидкости более, чем на порядок, превышает плотность газа. Потому приведенные рассуждения о недооценке плотности методом молекулярной динамики относятся к плотности жидкости.

В то же время значения поверхностного натяжения отличаются в ту же сторону. То есть заниженным значениям плотности, полученным методом молекулярной динамики, соответствует недооценка поверхностного натяжения для всех температур. Одной из целей работы является демонстрация возможностей метода молекулярной динамики в воспроизведении необходимых результатов, потому требуется найти способ установить связь между экспериментально получаемыми величинами и поверхностныи натяжением. 

Существуют способы получить совпадающие с экспериментальными значения поверхностного натяжения, например, вводя поправки к энергии взаимодействия или к самому значению поверхностного натяжения. В первом случае, если вводятся поправки к энергии, например, путем добавления хвостовых поправок, значения плотностей и энергий получаются некорректными в окрестности межфазной поверхности. В случае поправок к значению поверхностного натяжения, искажается значение парахора, что уменьшает возможность использования этих данных.

Зависимость поверхностного натяжения от разности плотностей представлена на рисунке~\cref{fig:parachor}. Формула~\cref{eq:parachor_definition} удовлетворительно описывает как поведение экспериментальной зависимости, так и полученной с помощью молекулярной динамики. Полученное для экспериментальной зависимости значение $P_0 = 156 \text{эрг}^{1/4} \text{см}^{5/2} /\text{г}$, в то время как значение, полученное с помощью молекулярной динамики $P_0 = 147 \pm \text{эрг}^{1/4} \text{см}^{5/2} /\text{г}$. В~\cite{escobedo} приводится значение $P_0 = 159 \text{эрг}^{1/4} \text{см}^{5/2} /\text{г}$. Хотя полученная зависимость не соответствует экспериментальной, полученное значение парахора лежит в пределах 10 процентов. Это позволяет предсказывать значения поверхностного натяжения для известных плотностей жидкой и газовой фаз, которые легко получить из эксперимента. Подобная ситуация наблюдалась для модели TraPPE-UA, когда модель не воспроизводит значения вязкости и диффузии для отдельных плотностей, но хорошо показывает зависимость вязкости от обратной диффузии.

\begin{figure}[ht]
    \centerfloat{
        \includegraphics[scale=0.75]{parachor.png}
    \caption{Зависимость поверхностного натяжения от разности плотностей этана в жидкой
и газовой фазах, полученные с помощью МД-моделирования чистого этана (квадраты) и в
эксперименте [2] (круги). Цифры показывают температуру в К, при которой
получена соответствующая точка.}
    }\label{fig:parachor}
\end{figure}

\subsection{Поверхностное натяжение многокомпонентных систем}\label{subsec:ch3/sect2/sub3}

Была получена изотерма 213~K смеси метана и этана для диапазона давлений от 4 до 40~атм. Полученные зависимости убывают с ростом доли метана в жидкой и газовой фазах, притом в случае жидкой фазы характер зависимости близок к линейному. Уменьшение поверхностного натяжения может быть следствием убывания разности плотностей при приближении к критической точке. Та же закономерность была отмечена в \cite{kaverin2006}. 

Линейной оказалась так же зависимость поверхностного натяжения от давления. Аналогично обсужденным выше, эта зависимость является убывающей с приближением давления к критической точке. Подобный характер зависимости был получен в \cite{kaverin2006}.

\begin{figure}[ht]
    \centerfloat{
        \hfill
        \subcaptionbox[List-of-Figures entry]{\label{fig:sigma_x}}{%
            \includegraphics[width=0.5\linewidth]{sigma_x.png}}
        \hfill
        \subcaptionbox{\label{fig:sigma_P}}{%
            \includegraphics[width=0.5\linewidth]{sigma_P.png}}
    }
    \caption{Зависимость поверхностного натяжения смеси метана и этана от доли метана в жидкой(треугольники) и газовой(круги) фазе(а) и от давления(б) для температуры 213 К. Цифры показывают давление соответствующей точки в атм.}\label{fig:sigma}
\end{figure}

Следовательно, независимо от числа компонентов в системе, основные зависимости не меняют вид. Это позволяет использовать данные молекулярного моделирования для получения наклона зависимости. Что, в свою очередь, совместно с известными из эксперимента плотностями, позволит получить корректную зависимость поверхностного натяжения смеси и от доли метана в жидкой фазе, и от давления.

\subsection{Выводы из третьей главы}

Методом молекулярной динамики рассчитано поверхностное натяжение как для чистого этана, так и для смеси метана и этана. Использована модель TraPPE-UA, которая хорошо описывает фазовые равновесия. Она не является полноатомной, что позволяет уменьшить время, которое требуется на расчеты. Объединение атома углерода с атомами водорода в псевдоатом сокращает время расчета примерно в 5 раз. Однако, из-за этого страдает точность.

Для чистого этана был исследован диапазон давлений 4-40~атм. Рассмотрены температуры от 203~К до 253~К. Расчеты для более высоких температур требовали большего времени на установление равновесия. Расчеты для низких температур не были бы подкреплены экспериментальными данными. 

По данным, полученным для чистого этана, видно, что происходит недооценка как плотности, так и поверхностного натяжения. Это связано с необходимостью отключения дальнодействующих поправок к потенциалу межатомного взаимодействия. В двухфазной системе имеется разница между плотностями фаз, что противоречит предположению о гомогенности системы, которое лежит в основе вывода формулы таких дальнодействующих поправок. Из этого следует, что использованный метод не подходит для прямого расчета поверхностного натяжения.

Рассчитанный парахор оказался в пределах 10 процентов от экспериментальной величины. Таким образом, метод молекулярной динамики дает возможность предсказывать характер зависимости поверхностного натяжения от разности давлений жидкой и газовой фазы. Следовательно, при наличии значения поверхностного натяжения из известной разности давлений можно определить поверхностное натяжение в некотором диапазоне разности давлений с помощью полученного методом МД значения парахора. 

Для смеси был проведён расчёт изотермы 213~К. Получена линейная зависимость поверхностного натяжения от давления, что совпадает с полученными из теоретических соображений результатами. 
